\subsection*{2.3. 有限性定理（等特征零情形）}

\subsubsection*{定理 2.3.1.}
\underline{设 $S$ 为 0.2.5 中的 Hensel 离散赋值环谱, $\Lambda$ 为 Noether 挠环, $f:X\longrightarrow S$ 为有限型态射. 假设 $S$ 为等特征零. 则对 $X_\eta$ 上的每个可构造 $\Lambda$-模层 $F$, $X_{\bar s}$ 上的层 $R^i\Psi_{\bar\eta}(F)$ 均可构造.}

先重复 SGA~4, XIV, 1.10 证明开头的论证. 应用 SGA~4, IX, 2.14(ii). 存在 $F$ 的消解 $F^\bullet$
\begin{center}
\begin{tikzcd}[column sep=large, row sep=large]
{} & F \arrow[d] & {} & {} \\
0 \arrow[r] & F^0 \arrow[r] & F^1 \arrow[r] & \cdots
\end{tikzcd}
\end{center}
其各项形如 $\bigoplus_i p'_{i*}G_i$, 其中 $p'_i:Y_{i\eta}\longrightarrow X_\eta$ 有限, $G_i$ 为可构造常值层. 可假设 $Y_{i\eta}$ 为约化概形; 设 $p_i:Y_i\longrightarrow X$ 为 $X$ 在 $Y_{i\eta}$ 中的正规化. 使用消解 $F^\bullet$ 的谱序列, 并对有限态射 $p_i$ 这一初等情形应用 (2.1.7.1), 可见只须对每一对 $(Y_i,G_i)$ 证明定理.

接下来, 对 $F$ 为 $\Lambda$-模 $M$ 所定义的常值层的情形证明定理. 回顾对某个适当的 $n$, $\Lambda$ 是 $\mathbb Z/n$-代数. 因此模 $M$ 可分解为 $\Lambda/\ell_i\Lambda$ 上的模 $M_i$, 其中每个 $\ell_i\mid n$ 为适当的素数. 这将我们归约到 $\Lambda$ 为 $\mathbb Z/\ell$-代数、$\ell$ 为素数的情形. 公式 (2.1.13.1) 此时给出
\[
M\otimes_{\mathbb Z/\ell}
R^i\Psi_{\bar\eta}(\mathbb Z/\ell)
\xrightarrow{\ \sim\ }
R^i\Psi_{\bar\eta}(M).
\]
于是只须处理常值层 $F=\mathbb Z/\ell$.

选择闭子集的下降序列
\[
X_\eta=X'_0\supset\cdots\supset X'_n
\]
使 $(X'_i-X'_{i+1})_{\mathrm{red}}$ 在 $\eta$ 上光滑, 并设 $X_i=\overline{X'_i}$ 为其在 $X$ 中的闭包. 常值层 $\mathbb Z/\ell$ 可分解为各层片 $X'_i-X'_{i+1}$ 上常值层 $\mathbb Z/\ell$ 的补零延拓. 因此:

\medskip
\noindent\underline{归约 1.}
只须对层 $j_!\mathbb Z/\ell$ 证明定理 2.3.1, 其中 $j:U\hookrightarrow X_\eta$ 为一个光滑开子集的包含映射.

由奇点解消, 存在态射 $g:X_1\longrightarrow X$, 使得

\medskip
\noindent(a) $g$ 固有, 并诱导同构 $g^{-1}(U)\xrightarrow{\sim}U$;

\noindent(b) $X_1$ 正则, 且 $X_1-g^{-1}(U)$ 为 $X_1$ 中的正规交叉除子, 是正则除子 $D_\alpha$ 的并.

设 $j_1:g^{-1}(U)\hookrightarrow X_1$. 则
\[
Rg_*j_{1!}\mathbb Z/\ell=j_!\mathbb Z/\ell.
\]
应用 (2.1.7.1) 与有限性定理 SGA~4, XIV, 1.1, 可见只须对 $(X_1,j_{1!}\mathbb Z/\ell)$ 证明定理 2.3.1.

有消解
\[
0\longrightarrow j_{1!}\mathbb Z/\ell
\longrightarrow\mathbb Z/\ell
\longrightarrow\bigoplus_\alpha(\mathbb Z/\ell)_{D_\alpha}
\longrightarrow
\bigoplus_{\alpha,\beta}
(\mathbb Z/\ell)_{D_\alpha\cap D_\beta}
\longrightarrow\cdots,
\]
从而:

\medskip
\noindent\underline{归约 2.}
只须对正则 $S$-概形 $X$ 上的常值层 $\mathbb Z/\ell$ 证明定理 2.3.1, 其中 $X_s$ 为 $X$ 中的正规交叉除子.

在最后这种情形中, 使用 I, 3.3 的显式计算; 该计算又使用纯性定理 SGA~4, XIX, 3.2 及特征零假设.
